\documentclass[11pt,a4paper]{article}
\usepackage[utf8]{inputenc}
\usepackage[vietnamese]{babel}
\usepackage{amsmath,amssymb}
\usepackage{geometry}
\geometry{a4paper, total={180mm,260mm}, left=15mm, right=15mm, top=18mm, bottom=18mm}

\usepackage{tikz}
\usetikzlibrary{shapes.geometric, arrows.meta, positioning, calc}
\tikzset{
  arr/.style={-{Stealth[scale=0.8]}},
  every node/.append style={align=center}
}
\usepackage{tcolorbox}
\tcbuselibrary{skins,breakable}
\usepackage{fontawesome}
\usepackage{fancyhdr}
\usepackage{enumitem}

\pagestyle{fancy}
\fancyhf{}
\renewcommand{\headrulewidth}{0.5pt}
\renewcommand{\footrulewidth}{0.5pt}
\fancyhead[L]{\small\bfseries\color{violet!85!black}\faBook\ SỔ TAY KIẾN THỨC TOÁN 2}
\fancyhead[R]{\small\bfseries\color{violet!85!black}Lớp Toán Thầy Lĩnh -- loptoanthaylinh.vn}
\fancyfoot[C]{\small\color{gray}-- Sổ tay kiến thức: Bài 43. Bảng chia 2 --}

\setlength{\parindent}{0pt}
\setlength{\parskip}{0pt}

\begin{document}

% --- BANNER BÀI HỌC ---
\begin{tcolorbox}[
  enhanced,
  colback=violet!8!white,
  colframe=violet!70!black,
  arc=3mm, boxrule=1.2pt,
  left=3mm, right=3mm, top=2mm, bottom=2mm,
  halign=center
]
  {\LARGE\bfseries\color{violet!90!black}\faBookmark\ BÀI 43: BẢNG CHIA 2}\\[1mm]
  {\itshape\footnotesize\color{black!75} Hình thành bảng chia 2 từ bảng nhân 2; kĩ năng tính nhẩm nhanh trên 1 trang A4}
\end{tcolorbox}

\vspace{-2mm}

\noindent
\begin{minipage}[t]{0.49\textwidth}\vspace{0pt}
\begin{tcolorbox}[
  enhanced,
  colback=violet!4!white,
  colframe=violet!75!black,
  arc=2mm, boxrule=0.8pt,
  left=2mm, right=2mm, top=1.8mm, bottom=1.8mm,
  title={\footnotesize\bfseries\color{white}\faKey\ 1. Bảng chia 2 chuẩn mực}
]
{\footnotesize
\begin{tabular}{l@{\quad}l}
  $2 : 2 = \mathbf{1}$ & $12 : 2 = \mathbf{6}$ \\
  $4 : 2 = \mathbf{2}$ & $14 : 2 = \mathbf{7}$ \\
  $6 : 2 = \mathbf{3}$ & $16 : 2 = \mathbf{8}$ \\
  $8 : 2 = \mathbf{4}$ & $18 : 2 = \mathbf{9}$ \\
  $10 : 2 = \mathbf{5}$ & $20 : 2 = \mathbf{10}$
\end{tabular}
\vspace{1mm}
\begin{itemize}\setlength{\itemsep}{0.5pt}\setlength{\leftskip}{-2.5mm}
  \item Số chia đều là $2$.
  \item Số bị chia là dãy số chẵn tăng dần: $2, 4, 6, \dots, 20$.
  \item Thương tăng dần từ $1$ đến $10$.
\end{itemize}
}
\end{tcolorbox}
\end{minipage}
\hfill
\begin{minipage}[t]{0.49\textwidth}\vspace{0pt}
\begin{tcolorbox}[
  enhanced,
  colback=teal!4!white,
  colframe=teal!75!black,
  arc=2mm, boxrule=0.8pt,
  left=2mm, right=2mm, top=1.8mm, bottom=1.8mm,
  title={\footnotesize\bfseries\color{white}\faSliders\ 2. Kĩ năng tính nhẩm dựa vào bảng nhân 2}
]
{\footnotesize
\begin{itemize}\setlength{\itemsep}{1pt}\setlength{\leftskip}{-2.5mm}
  \item Muốn tính $14 : 2 = ?$, ta nhẩm:\\
  $2$ nhân mấy bằng $14$?\\
  Vì $2 \times 7 = 14$ nên $14 : 2 = \mathbf{7}$.
  \item \textbf{Tìm một nửa (chia đôi):}\\
  Chia cho $2$ chính là tìm \textbf{một nửa} của số đó.\\
  Một nửa của $10$ là $10 : 2 = 5$.\\
  Một nửa của $18$ là $18 : 2 = 9$.
\end{itemize}
}
\end{tcolorbox}
\end{minipage}

\vspace{-1.5mm}
\begin{tcolorbox}[
  enhanced,
  colback=orange!4!white,
  colframe=orange!80!black,
  arc=2mm, boxrule=0.8pt,
  left=2mm, right=2mm, top=1.5mm, bottom=1.5mm,
  title={\footnotesize\bfseries\color{white}\faSitemap\ 3. Cầu nối 2 chiều: Bảng nhân 2 $\longleftrightarrow$ Bảng chia 2}
]
\centering
\begin{tikzpicture}[x=1.8cm, y=0.7cm]
  \node[rectangle, rounded corners=2pt, draw=violet!80!black, fill=violet!15, font=\scriptsize\bfseries] (n1) at (0,1) {$2 \times 4 = 8$};
  \node[rectangle, rounded corners=2pt, draw=teal!80!black, fill=teal!15, font=\scriptsize\bfseries] (c1) at (3,1) {$8 : 2 = 4$};
  \draw[<->, thick, orange!80!black] (n1) -- (c1);

  \node[rectangle, rounded corners=2pt, draw=violet!80!black, fill=violet!15, font=\scriptsize\bfseries] (n2) at (0,0) {$2 \times 8 = 16$};
  \node[rectangle, rounded corners=2pt, draw=teal!80!black, fill=teal!15, font=\scriptsize\bfseries] (c2) at (3,0) {$16 : 2 = 8$};
  \draw[<->, thick, orange!80!black] (n2) -- (c2);
\end{tikzpicture}
\end{tcolorbox}

\vspace{-1.5mm}
\noindent
\begin{minipage}[t]{0.49\textwidth}\vspace{0pt}
\begin{tcolorbox}[
  enhanced,
  colback=green!4!white,
  colframe=green!60!black,
  arc=2mm, boxrule=0.8pt,
  left=2mm, right=2mm, top=1.8mm, bottom=1.8mm,
  title={\footnotesize\bfseries\color{white}\faCheckCircle\ 4. Ví dụ thực hành chuẩn mực}
]
{\footnotesize
\textbf{1. Tính nhẩm:}\\
$12 : 2 = \mathbf{6}$; $18 : 2 = \mathbf{9}$; $10 : 2 = \mathbf{5}$.
\textbf{2. Bài toán lời văn:}\\
Có $16$ bông hoa cắm đều vào $2$ lọ. Hỏi mỗi lọ có bao nhiêu bông hoa?\\
\textit{Bài giải:}\\
Mỗi lọ có số bông hoa là:
$$16 : 2 = 8\text{ (bông hoa)}$$
Đáp số: $8$ bông hoa.
}
\end{tcolorbox}
\end{minipage}
\hfill
\begin{minipage}[t]{0.49\textwidth}\vspace{0pt}
\begin{tcolorbox}[
  enhanced,
  colback=red!4!white,
  colframe=red!75!black,
  arc=2mm, boxrule=0.8pt,
  left=2mm, right=2mm, top=1.8mm, bottom=1.8mm,
  title={\footnotesize\bfseries\color{white}\faExclamationTriangle\ 5. Mẹo tránh bẫy trắc nghiệm}
]
{\footnotesize
\begin{itemize}\setlength{\itemsep}{1pt}\setlength{\leftskip}{-2.5mm}
  \item \textbf{Chia cho 2 khác trừ đi 2:}\\
  $6 : 2 = 3$, nhưng $6 - 2 = 4$. Tuyệt đối không nhầm phép chia với phép trừ!
  \item Số bị chia trong bảng chia 2 luôn là số chẵn.
\end{itemize}
}
\end{tcolorbox}
\end{minipage}


\end{document}
